(1) 一维两束等强反向激光照射质量 \(m\) 的两能级原子，波数 \(k\)、自然线宽 \(\Gamma\)、红失谐 \(\delta<0\)，每束低速散射率 \(R(\Delta)=(\Gamma/2)s/[1+s+4\Delta^2/\Gamma^2]\)，其中 \(s\) 是每束有效饱和参数。考虑Doppler移频 \(\Delta_\pm=\delta\mp kv\)，求平均辐射压力 \(F(v)\)。
(2) 将 \(F(v)\) 展开到 \(v\) 一阶，求摩擦系数 \(\alpha\) 使 \(F\simeq-\alpha v\)，给出 \(\alpha>0\) 的失谐条件和小速度线性化条件。
(3) 采用简化的一维随机反冲模型：每次散射循环的动量方差为 \(2(\hbar k)^2\)，两束总循环率为 \(2R(\delta)\)。求Fokker–Planck约定 \(d\langle(\Delta p)^2\rangle/dt=2D_p\) 下的扩散系数与稳态 \(k_BT=D_p/\alpha\)。
(4) 在 \(s\ll1\) 下优化失谐得到Doppler最低温度，检查 \(|\delta|\to0,\infty\) 的升温极限，并说明实际三维自发辐射角分布、亚Doppler机制与反冲温度对本模型的修正。
